% Imporditud: tools/import_eksperimendid.py
% Allikas: Eesti füüsikaolümpiaadi eksperimendi ülesannete
%   kogu (Rael Kalda, 2023), ülesanne Ü128, lahendus L128.
\setAuthor{EFO žürii}
\setRound{lõppvoor}
\setYear{2002}
\setNumber{G E1}
\setDifficulty{4}
\setTopic{geomeetriline optika}
\setVahendid{nõguslääts, joonlaud, ruuduline paber}
\setPunktid{12}
\setAllikas{Eksperimendi kogu Ü128, lahendus lk 99}
\setLahendusseis{toores}

\prob{Läätse optiline tugevus}
Hinnata läätse optilist tugevust. Põhjenduseks esitada ka joonis ning katse kirjeldus.

\hint

\solu
Tõmbame pliiatsiga paberile lõigu pikkusega l, näiteks l = 20 cm. Viime silma paberist kaugusele L, näiteks L = 30 cm. Viime läätse paberist sellise kaugusele a, et joonistatud lõik paistab läbi läätse täies pikkuses ulatudes servast servani. Teeme paberile sobivas mõõtkavas täpse joonise. Juuresoleval pildil OA = l/2, P B = d/2 (d on läätse diameeter), OC = L ja OP = a. Konstrueerime punkti F fokaaltasandil nii, nagu näidatud joonisel. Mõõdame fookuskauguse f , optiline tugevus D = 1/f .

A

F B

O Q P C

f

Võib ka avaldada analüütiliselt. Selleks paneme kirja sarnaste kolmnurkade sar-

nasustingimused:

CQ CP PQ PO =, = . F Q P B F Q AO $-$ P B

Jagades nende võrduste vasakud ja paremad pooled saame välja taandada pikkuse F Q. Saadud võrrandist on lihtne avaldada D = 1/f , kui arvestame, et

PQ = f, CQ = f + L $-$ a, CP = L $-$ a, P O = a, P B = d/2, AO = l/2.

() 1l 1 D= d $-$1 $-$ L $-$ a $\approx$ 8 dptr. a
\probend
